This problem explores some aspects of the two main methods of exoplanet
detection: radial velocity and transit. Throughout this problem we shall
consider a particular system of a single planet (P) in a circular orbit with
radius $a$ around a solar-type star (S). We shall refer to this system as the
``SP system''.

\medskip\noindent\textbf{Stellar parameters}

\begin{description}
\item[(D01.1)] The V-band apparent magnitude of the star S is
  $(7.65\pm0.03)$\,mag, the parallax is $(20.67\pm0.05)$\,milliarcsecond and
  the bolometric correction (BC) is $-0.0650$\,mag. Thus the star has a higher
  bolometric luminosity than its luminosity in the V-band.

  Estimate the mass of the star, $M_{\mathrm{s}}$ (in units of $M_\odot$),
  assuming a mass-luminosity ($M$--$L$) relation of the form $L \propto M^4$.
  Also estimate the uncertainty in $M_{\mathrm{s}}$. You may need
  $\frac{d}{dx}\ln x = \frac{1}{x}$. \hfill[8]
\end{description}

\medskip\noindent\textbf{Radial Velocity method}

The radial velocity method uses the Doppler shift
$\delta\lambda \equiv \lambda_{\mathrm{obs}} - \lambda_0$ between the observed
wavelength $\lambda_{\mathrm{obs}}$ and the rest wavelength $\lambda_0$ of a
known spectral line to detect an exoplanet and determine its characteristics.

The figure below shows the $\delta\lambda$ for the Fe\,I line
($\lambda_0 = 543.45\times10^{-9}$\,m) as a function of time as observed for
the SP system.

\ioaafig{2025_DA_D01-f1.pdf}
% f1 = scatter plot of delta-lambda (in 10^-15 m) against time (0-24 days)

The radial velocity semi-amplitude $K$ is defined as
$K \equiv (v_{\mathrm{r,max}} - v_{\mathrm{r,min}})/2$ where
$v_{\mathrm{r,max}}$ and $v_{\mathrm{r,min}}$ are the minimum and maximum
radial velocities, respectively.
% sic: the original indeed says "the minimum and maximum radial velocities"
% in this order for v_r,max and v_r,min
For a circular planetary orbit the semi-amplitude $K$ can be written as:
\[
K = \left(\frac{2\pi G}{T}\right)^{1/3}
    \frac{M_{\mathrm{p}}\sin i}{(M_{\mathrm{p}} + M_{\mathrm{s}})^{2/3}}
\]
where $T$ is the period, $i$ is the inclination of the planetary orbit (angle
between the normal to the orbital plane of the planet and the line of sight of
the observer), $M_{\mathrm{p}}$ and $M_{\mathrm{s}}$ are the masses of the
planet and the star, respectively.

\begin{description}
\item[(D01.2)] Use the graph given in the Summary Answersheet to answer the
  following.
\item[(D01.2a)] Draw a smooth curve associated with the observed data shown in
  the graph. \hfill[2]
\item[(D01.2b)] Select appropriate points on your drawn curve and use suitable
  methods to determine $T$ and $K$ along with respective uncertainties. All
  data points used for the calculation of $T$ and $K$ must be shown in the
  table in the Summary Answersheet. Use the rest of the table to show your
  intermediate calculations, as needed, with appropriate headers. \hfill[11]
\item[(D01.2c)] Find the minimum mass of the planet $M_{\mathrm{p,min}}$ (in
  $M_\odot$), and its corresponding uncertainty, assuming
  $M_{\mathrm{p}} \ll M_{\mathrm{s}}$. \hfill[5]
\item[(D01.2d)] Using the value of $M_{\mathrm{p,min}}$ estimated in part
  (D01.2c), calculate the minimum value of the semi-major axis of the planet's
  orbit, $a_{\mathrm{min}}$, in au and its uncertainty. \hfill[4]
\end{description}

\medskip\noindent\textbf{Transit method (without limb darkening)}

The schematic diagram of a planet transit (not drawn to scale) is shown
below. Initially, we shall assume the stellar disk to have a uniform average
intensity with some intrinsic noise due to the star itself.

\ioaafig{2025_DA_D01-f2.pdf}
% f2 = schematic: face-on and side-on views of the transit (uniform disk),
% with R_s, bR_s, R_p, inclination i, and the schematic lightcurve showing
% Delta, t_F and t_T

The lightcurve of the normalized intensity, $I$, as a function of time $t$ is
shown in the schematic diagram of the transit above. The average stellar
intensity outside the transit is taken as unity. The maximum decrease in the
intensity is given by $\Delta$ in the normalized light curve. For a uniformly
bright stellar disk, the radius of the planet, $R_{\mathrm{p}}$, is related to
$\Delta$ as
\[
\left(\frac{R_{\mathrm{p}}}{R_{\mathrm{s}}}\right)^2 = \Delta,
\]
where $R_{\mathrm{s}}$ is the radius of the star.

The total duration of transit (when part or all of the planet covers the
stellar disk) is given by $t_{\mathrm{T}}$, while $t_{\mathrm{F}}$ gives the
duration when the planet is fully in front of the stellar disk. The ``impact
parameter'' $b$ is the projected distance between the planet and centre of the
stellar disk at the mid-point of the transit, in units of the stellar radius,
$R_{\mathrm{s}}$.

For a nearly edge-on star-planet orbit, the impact parameter is given by the
formula
\[
b = \left[\frac{(1-\sqrt{\Delta})^2
      - (t_{\mathrm{F}}/t_{\mathrm{T}})^2(1+\sqrt{\Delta})^2}
     {1 - (t_{\mathrm{F}}/t_{\mathrm{T}})^2}\right]^{1/2}
\]

\begin{description}
\item[(D01.3)] For the SP system, the stellar radius is known to be
  $R_{\mathrm{s}} = 1.20\,R_\odot$, and the transit of the planet is indeed
  visible. Using the minimum orbital radius, $a_{\mathrm{min}}$, estimated in
  part (D01.2d), find the minimum value, $i_{\mathrm{min}}$, of the
  inclination angle. \hfill[3]
\end{description}

Assuming a stellar disk of uniform brightness, the transit lightcurve would
look like as shown below.

\ioaafig{2025_DA_D01-f3.pdf}
% f3 = transit lightcurve: normalized intensity (0.980-1.000) against days
% relative to transit midpoint (-0.15 to +0.15)

\begin{description}
\item[(D01.4)] Using the given lightcurve answer the following questions. For
  your reference the above lightcurve is also given in the Summary
  Answersheet.
\item[(D01.4a)] Estimate the values of $t_{\mathrm{T}}$ and $t_{\mathrm{F}}$
  in days by marking appropriate readings on the graph. \hfill[3]
\item[(D01.4b)] Estimate the mean value of $\Delta$, by marking appropriate
  readings on the graph and hence find $R_{\mathrm{p}}$ in units of
  $R_\odot$. \hfill[2]
\item[(D01.4c)] Determine the value of $i$ in degrees assuming the orbital
  radius to be $a_{\mathrm{min}}$. \hfill[2]
\end{description}

\medskip\noindent\textbf{Introducing limb darkening}

So far we have assumed the stellar disk to be uniformly bright. In reality,
the observed brightness of the stellar disk is not uniform due to ``limb
darkening'' --- an optical effect where the central part of the stellar disk
appears brighter than the edge, or the ``limb''.

The limb darkening effect can be measured by the relative intensity
$J(\theta) \equiv \frac{I(\theta)}{I(0)}$, where $\theta$ is the angle between
the normal to the stellar surface at a point and the line joining the observer
to that point, $I(\theta)$ is the observed intensity of the stellar disk at
that point ($I(0)$ being the intensity at the centre of the stellar disk). For
a distant observer, $\theta$ varies from $\theta = 0$ (centre of the disk) to
$\theta \approx 90^\circ$ (edge of the disk).

\begin{description}
\item[(D01.5)] The table below gives measured $J(\theta)$ at a certain
  wavelength for the Sun. We shall assume that the same limb darkening profile
  holds for the star S.

  \begin{center}
  \begin{tabular}{cc}
  \toprule
  $\theta$ & $J(\theta)$\\
  \midrule
  $0^\circ$  & 1.000\\
  $10^\circ$ & 0.994\\
  $15^\circ$ & 0.984\\
  $20^\circ$ & 0.971\\
  $25^\circ$ & 0.950\\
  $30^\circ$ & 0.943\\
  $40^\circ$ & 0.883\\
  $50^\circ$ & 0.794\\
  $60^\circ$ & 0.724\\
  $70^\circ$ & 0.595\\
  $80^\circ$ & 0.475\\
  $90^\circ$ & 0.312\\
  \bottomrule
  \end{tabular}
  \end{center}

  The limb darkening profile can be modelled by a quadratic formula:
  \[
  J(\theta) = 1 - a_1(1-\cos\theta) - a_2(1-\cos\theta)^2,
  \]
  where $a_1$ and $a_2$ are two constants.

  We shall estimate the unknown coefficients $a_1$ and $a_2$ from the given
  data by making a plot with suitable variables.
\item[(D01.5a)] Choose a pair of variables $(x_1, y_1)$ which are suitable
  functions of $\theta$ and $J$, that you want to plot along $x$ and $y$ axes,
  respectively, to determine $a_1$ and $a_2$. Write the expressions for $x_1$
  and $y_1$.

  If you need to define additional variables for additional plots, define them
  as $(x_2, y_2)$, etc. \hfill[2]
\item[(D01.5b)] Tabulate the values necessary for your plots. \hfill[4]
\item[(D01.5c)] Plot the newly defined variables on the given graph paper
  (mark your graph as ``D01.5c''). \hfill[7]
\item[(D01.5d)] Obtain $a_1$ and $a_2$ from the plot. Uncertainties on the
  values are not needed. \hfill[7]
\end{description}

\medskip\noindent\textbf{Transit in the presence of limb darkening}

Now, we consider planetary transits across a limb darkened stellar disk. In
the presence of limb darkening, which we shall model by the quadratic formula
of $J(\theta)$ given above, the average observed intensity of the entire
stellar disk (without any transit), $\langle I\rangle$, is given by:
\[
\langle I\rangle = \left(1 - \frac{a_1}{3} - \frac{a_2}{6}\right) I(0)
\]
Further, the dip in the light caused by the transiting planet now depends not
only on the relative size of the planet and the star,
$\left(\frac{R_{\mathrm{p}}}{R_{\mathrm{s}}}\right)$, but also on the
intensity profile of the stellar disk along the transit chord, which in turn,
depends on the angle of inclination, $i$.

The schematic diagram below (not drawn to scale) shows the configuration. Note
that the brighter part of the star is shown in a darker shade, while the
planet is shown as a black dot.

\ioaafig{2025_DA_D01-f4.pdf}
% f4 = schematic: face-on view (transit chord through point C) and side-on
% view (angle theta_C between line of sight and surface normal) of a transit
% across a limb-darkened disk

Here the relation between
$\left(\frac{R_{\mathrm{p}}}{R_{\mathrm{s}}}\right)$ and the measured $\Delta$
from the light curve is
\[
\Delta = \frac{I(\theta_{\mathrm{C}})}{\langle I\rangle}
  \left(\frac{R_{\mathrm{p}}}{R_{\mathrm{s}}}\right)^2,
\]
where $I(\theta_{\mathrm{C}})$ is the intensity of the stellar disk at the
midpoint of the transit chord (point C in the figure above),
$\theta_{\mathrm{C}}$ being the angle between the line of sight and the normal
to the surface at that point. From the above it is obvious that for a given
star, the same value of $\Delta$ can be produced by many combinations of the
planet size, $R_{\mathrm{p}}$, and the inclination angle $i$.

\begin{description}
\item[(D01.6)] It is possible to uniquely determine both $R_{\mathrm{p}}$ and
  $i$ by using data from transit lightcurves at two wavelengths, say,
  $\lambda_{\mathrm{B}}$ (blue) and $\lambda_{\mathrm{R}}$ (red). The limb
  darkening coefficients for these two wavelengths are given below:

  \begin{center}
  \begin{tabular}{ccc}
  \toprule
  Wavelength & $a_1$ & $a_2$\\
  \midrule
  $\lambda_{\mathrm{B}}$ & 0.82 & 0.05\\
  $\lambda_{\mathrm{R}}$ & 0.24 & 0.20\\
  \bottomrule
  \end{tabular}
  \end{center}
\item[(D01.6a)] Choose the correct statement among the following that
  describes the relation between the maximum depth of the transit ($\Delta$)
  for $\lambda_{\mathrm{B}}$ and the inclination angle ($i$) of the orbit and
  tick ($\checkmark$) it in the Summary Answersheet.

  \begin{itemize}
  \item[(A)] $\Delta$ increases with decreasing $i$.
  \item[(B)] $\Delta$ decreases with decreasing $i$.
  \item[(C)] $\Delta$ is independent of $i$.
  \end{itemize}
  \hfill[2]
\item[(D01.6b)] The maximum depth of the transit ($\Delta$) for the ``SP
  system'' was measured to be 0.0182 and 0.0159 for $\lambda_{\mathrm{B}}$ and
  $\lambda_{\mathrm{R}}$, respectively.

  Draw schematic transit light curves for both $\lambda_{\mathrm{B}}$ and
  $\lambda_{\mathrm{R}}$ on the given grid and label the curves by ``B'' and
  ``R'' respectively. Assume that the total transit duration is same for both
  wavelengths. The curves need not be to scale, but should represent the
  shapes of the light curves correctly. \hfill[4]
\item[(D01.7)] We shall use a graphical method to find the values of
  $R_{\mathrm{p}}$ and $i$ for the SP system using measurements of $\Delta$ at
  $\lambda_{\mathrm{B}}$ and $\lambda_{\mathrm{R}}$.
\item[(D01.7a)] Write an appropriate expression connecting the relevant
  variables that are to be plotted. (Hint: You may consider $i$ or $b$, and
  $R_{\mathrm{p}}$ among the relevant variables.) \hfill[6]
\item[(D01.7b)] Tabulate the appropriate quantities that are to be
  plotted. \hfill[5]
\item[(D01.7c)] Draw a suitable graph and mark it as ``D01.7c''. \hfill[7]
\item[(D01.7d)] Estimate the values of $R_{\mathrm{p}}$ (in $R_\odot$) and $i$
  (in degrees) from the graph. \hfill[4]
\item[(D01.8)] Based on the results obtained in this problem, indicate whether
  the planet P is ``ROCKY'' or ``GASEOUS'' by ticking ($\checkmark$) the
  appropriate box in the Summary Answersheet. \hfill[2]
\end{description}
